\label{sec:resumo_atividades}

\mbox{Usaram-se} os eventos simulados da etapa anterior, cujos arquivos datam de outubro de 2025: depósitos de \qty{100}{\MeV}, reescalados para \qty{10}{\MeV}, num octante do terceiro módulo do detector distante (FD3) do DUNE (\textit{Deep Underground Neutrino Experiment}), lidos por \num{6912} canais. Em dezembro de 2025, \mbox{testou-se} um protótipo de rede neural em grafos (GNN), seguido, em fevereiro e março de 2026, dos codificadores GATv2 e GPS, com cabeça de regressão hierárquica. Em junho, o código foi reconstruído como pacote instalável; corrigiram\discretionary{-}{-}{-}se dois defeitos de geometria e três canais quentes, os eventos de treino foram espelhados nos outros octantes e os grafos foram restritos aos canais com luz, com \num{17} atributos por nó e \num{23} por aresta. Em julho, um modelo simples de \textit{gradient boosting} (árvores de decisão ajustadas em sequência) superou a GNN da época, e os \textit{boosters} XGBoost e LightGBM entraram como referência.

Em 4 de agosto de 2026, \mbox{ajustaram-se} os \textit{boosters}, e uma primeira comparação treinou \num{11} dos \num{14} codificadores com um quarto dos eventos de treino. Um diagnóstico da diferença para os \textit{boosters} levou ao uso dos \num{7490} eventos de treino e \num{33} escalares do evento, que resumem sua luz. Na comparação final, de 8 de agosto, \num{11} codificadores de capacidade equiparada, com desvio máximo de \qty{1.6}{\percent} em relação a \num{1843395} parâmetros, somaram \qty{11.42}{\hour} de treino.

Neste relatório, sem novo treinamento, os \textit{boosters}, ajustados e treinados em luz com aumento de dados, foram reavaliados na mesma luz dos codificadores, e \mbox{calcularam-se} os intervalos de confiança e o baricentro da luz, nos mesmos eventos de teste. Os produtos são um artigo com os resultados para o Encontro de Iniciação Científica do Instituto Tecnológico de Aeronáutica (ENCITA) de 2026, \textit{slides} do projeto e um repositório de \num{155} \textit{commits} (versões registradas) e \num{172} arquivos Python, com \num{18025} linhas e \num{130} funções de teste.

Nos \num{937} eventos de teste, os quatro codificadores de menor erro (duas variantes do GATv2 e duas do GPS), indistinguíveis entre si, têm erro tridimensional médio (distância média do vértice previsto ao verdadeiro) de \num{25.95} a \qty{26.27}{\cm} e raiz do erro quadrático médio dos três eixos juntos (RMSE combinado) de \num{18.55} a \qty{18.80}{\cm}. O XGBoost, na mesma luz, tem RMSE combinado de \qty{20.95}{\cm}, mas erra menos nos eventos mais claros. A rede da etapa anterior tem \qty{20.86}{\cm}, valor calculado a partir dos RMSE publicados por eixo, de outra partição \cite{vicentin2024study}. O APEX (\textit{Aluminum Profiles with Embedded X-ARAPUCA}), de outra simulação \cite{marinho2025apex}, tem largura combinada (média quadrática das larguras de resolução dos três eixos, lidas de um gráfico) de \qty{38.0}{\cm} na luz média por evento do teste com o seu rendimento luminoso médio, \qty{31.8}{\cm} com o mínimo e \qty{25.9}{\cm} em energia igual. Em $z$, o eixo longo, contra \num{20.15} a \qty{20.46}{\cm} dos quatro codificadores, o APEX tem \qty{24.8}{\cm} com o rendimento médio, acima do intervalo de confiança de cada um, e \qty{19.4}{\cm} com o mínimo, dentro dele; só em energia igual, em que detecta \num{3.4} vezes mais luz, é melhor (\qty{13.7}{\cm}).
